\documentclass[11pt,a4paper]{article}
\usepackage{layout}
\begin{document}

% 01 — Original 2026 cover; only OSDGs labels replaced with GOSD.
\begin{canvas}
\pdfbookmark[0]{封面}{cover}
\node[anchor=north west,inner sep=0pt,outer sep=0pt] at (current page.north west){\includegraphics[width=210mm,height=297mm]{.build/assets/cover2026-gosd.pdf}};
\end{canvas}

% 02 — Editorial board. Visual reading order preserved from the 2025 source.
\begin{canvas}
\pdfbookmark[0]{编委会}{board}
\Header{编委会}{teal}\Footer{02}
\Kicker{teal}{EDITORIAL BOARD}\Title{编委会}
\Subhead{72}{teal}{主编}\Rule{43}{77}{188}{line}
\Box{22}{85}{166}{\sffamily\fontsize{12}{20}\selectfont 王\quad 伟\hspace{13mm}李建忠}
\Subhead{108}{teal}{联合主席}\Rule{49}{113}{188}{line}
\Box{22}{121}{166}{\sffamily\fontsize{12}{20}\selectfont 王建民\hspace{13mm}张延川}
\Subhead{144}{teal}{副主席}\Rule{47}{149}{188}{line}
\Box{22}{157}{166}{\sffamily\fontsize{11}{19}\selectfont 陈\quad 钟\hspace{6mm}胡春明\hspace{6mm}黄\quad 韬\hspace{6mm}高\quad 跃\hspace{6mm}堵俊平\hspace{6mm}蒋\quad 涛}
\Subhead{181}{teal}{编委}\Rule{43}{186}{188}{line}
\Box{22}{196}{166}{\sffamily\fontsize{9.8}{21.5}\selectfont\setlength{\tabcolsep}{0pt}\begin{tabular*}{166mm}{@{\extracolsep{\fill}}llllllllll@{}}
吴发国&崔鹤鸣&于仕琪&周明辉&胡\ 钰&黄向东&王\ 博&余\ 跃&彭煜玮&陈\ 阳\\
顾\ 荣&金耀辉&游\ 兰&王\ 莹&张宇霞&何海洋&杨丽蕴&张俊霞&郭\ 雪&任旭东\\
许哲平&姜\ 宁&边思康&马红伟&郭\ 炜&谭中意&林旅强&江\ 波&陈\ 阳&陈\ 绪\\
贾\ 开&郑宏威&竺彩华&刘\ 广&隆云滔&乔嘉林&吴咏炜&瞿晶晶&于晓轶&邹\ 慧\\
夏小雅&赵生宇&武剑凌&方心雨&王\ 哲&汤文泰&董\ 超&彭佳恒&张\ 震&李沫南
\end{tabular*}}
\Source{263}{编委会沿用 2025 版，职务、名单及排列顺序按原版保留。}
\end{canvas}

% 03 — Contents I
\begin{canvas}
\pdfbookmark[0]{目录}{contents}
\Header{目录 · 上}{blue}\Footer{03}
\Kicker{blue}{CONTENTS / 01}\Title{目录}
\Box{22}{64}{166}{\sffamily\fontsize{10}{16}\selectfont\hyperlink{board.0}{编委会\quad 02}\hfill\hyperlink{summary}{执行摘要\quad 05}\hfill\hyperlink{intro}{引言：开源共创智能未来\quad 07}}
\Chapter{86}{01}{teal}{从开放协作到公共价值\hfill{\fontsize{10}{12}\selectfont 10}}
\Tocline{100}{1.1}{人工智能时代的开源协作：双向重塑\hfill 11}
\Tocline{108}{1.2}{开源发展目标：从开放资源到公共能力\hfill 14}
\Tocline{116}{1.3}{观察方法：跨平台的全球开源全景\hfill 16}
\Chapter{134}{02}{blue}{全球开源生态观测地图}
\Tocline{148}{2.1}{全球开源活跃地图}
\Tocline{156}{2.2}{全球开源贡献地图}
\Tocline{164}{2.3}{全球开源项目影响力地图}
\Tocline{172}{2.4}{全球开源人才地图}
\Tocline{180}{2.5}{全球开源资产地图}
\Tocline{188}{2.6}{全球开源多样性地图}
\Tocline{196}{2.7}{全球开源治理与社会公益地图}
\Chapter{207}{03}{gold}{全球开源影响力排行榜}
\Tocline{221}{3.1}{全球行政区划影响力排行榜}
\Tocline{229}{3.2}{全球开源企业影响力排行榜}
\Tocline{237}{3.3}{全球开源基金会影响力排行榜}
\Tocline{245}{3.4}{全球开源项目影响力排行榜}
\Tocline{253}{3.5}{全球新势力开源仓库与项目排行榜}
\Tocline{261}{3.6}{全球开源模型影响力排行榜}
\end{canvas}

% 04 — Contents II
\begin{canvas}
\Header{目录 · 下}{blue}\Footer{04}
\Kicker{blue}{CONTENTS / 02}\Title{由生态格局，走向共同能力}
\Chapter{68}{04}{teal}{全球开源技术全景洞察}
\Tocline{82}{4.1}{Agentic AI 全景洞察}
\Tocline{90}{4.2}{AI 大模型全景洞察}
\Tocline{98}{4.3}{云基础设施全景洞察}
\Tocline{106}{4.4}{数据库系统全景洞察}
\Tocline{114}{4.5}{RISC-V 全景洞察}
\Tocline{122}{4.6}{PyTorch 生态观察样本}
\Chapter{140}{05}{gold}{全球开源大模型发展洞察}
\Tocline{154}{5.1}{供给与贡献}
\Tocline{162}{5.2}{网络与影响力}
\Tocline{170}{5.3}{开放条件与分级}
\Chapter{188}{06}{violet}{未来展望与行动建议}
\Tocline{202}{6.1}{十二项关键洞察}
\Tocline{210}{6.2}{十大趋势判断与前瞻}
\Tocline{218}{6.3}{十大战略建议}
\Tocline{226}{6.4}{结语：开源共创智能未来}
\Box{22}{251}{166}{\sffamily\fontsize{10.2}{15}\selectfont\textbf{附录 A}\quad 指标体系与数据来源\hfill\textbf{致谢}}
\Source{265}{本次评审：前置内容 01—09 页，第一章 10—20 页；后续章节页码在合稿时生成。}
\end{canvas}

% 05 — Growth evidence and report distinctions
\begin{canvas}
\hypertarget{summary}{}\pdfbookmark[0]{执行摘要}{summary-bm}
\Header{执行摘要 · 数据与特色}{blue}\Footer{05}
\Kicker{blue}{EXECUTIVE SUMMARY}\Title{执行摘要}
\Subhead{59}{blue}{创造越快，越要看见全球协作与体系能力}
\Box{22}{70}{166}{\fontsize{10.3}{16.4}\selectfont 人工智能时代，代码、模型与数据加速增长。代码协作平台 GitHub 与模型资源平台 Hugging Face（HF）记录了这一变化。GOSD 从中国出发，以全球视野贯通“参与、供给、复用、治理与公共价值”，观察开放资源如何形成产业优势与公共能力。}
\Box{22.00}{93.00}{78.00}{\sffamily\bfseries\fontsize{11.5}{16}\selectfont\color{blue}GitHub 开发活动}
\Box{22.00}{102.00}{78.00}{\sffamily\fontsize{7.7}{11}\selectfont\color{muted}浅色：2025 · 深色：2026（十亿次）}
\Box{22.00}{113.00}{53.00}{\sffamily\bfseries\fontsize{9.4}{13}\selectfont 代码提交}
\Box{75.00}{113.00}{25.00}{\raggedleft\sffamily\bfseries\fontsize{9.6}{13}\selectfont\color{blue}约 +166\%}
\Rect{22}{120}{18.429}{3.4}{blue!25}
\Box{42.43}{119.00}{24.00}{\digits\fontsize{8.8}{12}\selectfont 12.9}
\Rect{22}{127}{49.000}{3.4}{blue}
\Box{73.00}{126.00}{24.00}{\digits\fontsize{8.8}{12}\selectfont 34.3}
\Box{22.00}{135.00}{53.00}{\sffamily\bfseries\fontsize{9.4}{13}\selectfont 自动化作业}
\Box{75.00}{135.00}{25.00}{\raggedleft\sffamily\bfseries\fontsize{9.6}{13}\selectfont\color{blue}约 +224\%}
\Rect{22}{142}{15.110}{3.4}{blue!25}
\Box{39.11}{141.00}{24.00}{\digits\fontsize{8.8}{12}\selectfont 14.0}
\Rect{22}{149}{49.000}{3.4}{blue}
\Box{73.00}{148.00}{24.00}{\digits\fontsize{8.8}{12}\selectfont 45.4}
\Box{22.00}{157.00}{53.00}{\sffamily\bfseries\fontsize{9.4}{13}\selectfont 合入修改提案}
\Box{75.00}{157.00}{25.00}{\raggedleft\sffamily\bfseries\fontsize{9.6}{13}\selectfont\color{blue}约 +142\%}
\Rect{22}{164}{20.278}{3.4}{blue!25}
\Box{44.28}{163.00}{24.00}{\digits\fontsize{8.8}{12}\selectfont 0.538}
\Rect{22}{171}{49.000}{3.4}{blue}
\Box{73.00}{170.00}{24.00}{\digits\fontsize{8.8}{12}\selectfont 1.30}
\Box{110.00}{93.00}{78.00}{\sffamily\bfseries\fontsize{11.5}{16}\selectfont\color{teal}HF 开放资产}
\Box{110.00}{102.00}{78.00}{\sffamily\fontsize{7.7}{11}\selectfont\color{muted}浅色：期初 · 深色：期末（万个）}
\Box{110.00}{113.00}{53.00}{\sffamily\bfseries\fontsize{9.4}{13}\selectfont 公共模型仓库}
\Box{163.00}{113.00}{25.00}{\raggedleft\sffamily\bfseries\fontsize{9.6}{13}\selectfont\color{teal}约 +22\%}
\Rect{110}{120}{40.226}{3.4}{teal!25}
\Box{152.23}{119.00}{24.00}{\digits\fontsize{8.8}{12}\selectfont 243}
\Rect{110}{127}{49.000}{3.4}{teal}
\Box{161.00}{126.00}{24.00}{\digits\fontsize{8.8}{12}\selectfont 296}
\Box{110.00}{135.00}{53.00}{\sffamily\bfseries\fontsize{9.4}{13}\selectfont 数据集}
\Box{163.00}{135.00}{25.00}{\raggedleft\sffamily\bfseries\fontsize{9.6}{13}\selectfont\color{teal}约 +41\%}
\Rect{110}{142}{34.839}{3.4}{teal!25}
\Box{146.84}{141.00}{24.00}{\digits\fontsize{8.8}{12}\selectfont 71.1}
\Rect{110}{149}{49.000}{3.4}{teal}
\Box{161.00}{148.00}{24.00}{\digits\fontsize{8.8}{12}\selectfont 100}
\Box{110.00}{157.00}{53.00}{\sffamily\bfseries\fontsize{9.4}{13}\selectfont 应用空间}
\Box{163.00}{157.00}{25.00}{\raggedleft\sffamily\bfseries\fontsize{9.6}{13}\selectfont\color{teal}约 +44\%}
\Rect{110}{164}{34.028}{3.4}{teal!25}
\Box{146.03}{163.00}{24.00}{\digits\fontsize{8.8}{12}\selectfont 100}
\Rect{110}{171}{49.000}{3.4}{teal}
\Box{161.00}{170.00}{24.00}{\digits\fontsize{8.8}{12}\selectfont 144}

\Box{22}{178}{166}{\sffamily\fontsize{7.8}{11.5}\selectfont\color{muted}图摘 1\quad 开发活动与开放资产增长；各指标独立标度，增幅按所列数值估算。}
\Box{22}{185}{166}{\sffamily\fontsize{7.8}{11.5}\selectfont\color{muted}来源：\href{\GH}{GitHub Octoverse 2026}；\href{\HF}{Hugging Face Summer 2026}。GitHub 各年为截至 9 月 30 日的 12 个月，含公开与私有仓库；HF 为 2026 年 1—8 月公共资源存量，非独立原创模型或新增产出。两者均与本报告主统计期不同。}
\Subhead{204}{teal}{本报告的特色}
\shade[left color=teal!10,right color=teal!2,rounded corners=2mm,draw=teal!17,line width=.45pt] \Pos{22}{215} rectangle \Pos{102}{242};
\path[fill=teal!9] \Pos{50.8}{242} .. controls \Pos{76.4}{241} and \Pos{82.8}{237} .. \Pos{101}{238} -- \Pos{101}{241} -- \Pos{50.8}{242} -- cycle;
\Box{26.00}{217.50}{72.00}{\sffamily\bfseries\fontsize{10.4}{15}\selectfont\color{teal}AI 全景 · 体系化观察}
\Box{26.00}{225.50}{72.00}{\fontsize{9.3}{14.1}\selectfont 以地图、排行、技术全景与模型网络，贯通参与、供给、复用与治理。}
\shade[left color=teal!10,right color=teal!2,rounded corners=2mm,draw=teal!17,line width=.45pt] \Pos{108}{215} rectangle \Pos{188}{242};
\path[fill=teal!9] \Pos{136.8}{242} .. controls \Pos{162.4}{241} and \Pos{168.8}{237} .. \Pos{187}{238} -- \Pos{187}{241} -- \Pos{136.8}{242} -- cycle;
\Box{112.00}{217.50}{72.00}{\sffamily\bfseries\fontsize{10.4}{15}\selectfont\color{teal}跨平台 · 全域数据视野}
\Box{112.00}{225.50}{72.00}{\fontsize{9.3}{14.1}\selectfont 联读 GitHub、HF、Gitee 等平台，从资源数量深入协作与复用关系。}
\shade[left color=teal!10,right color=teal!2,rounded corners=2mm,draw=teal!17,line width=.45pt] \Pos{22}{246} rectangle \Pos{102}{273};
\path[fill=teal!9] \Pos{50.8}{273} .. controls \Pos{76.4}{272} and \Pos{82.8}{268} .. \Pos{101}{269} -- \Pos{101}{272} -- \Pos{50.8}{273} -- cycle;
\Box{26.00}{248.50}{72.00}{\sffamily\bfseries\fontsize{10.4}{15}\selectfont\color{teal}中国出发 · 全球视野}
\Box{26.00}{256.50}{72.00}{\fontsize{9.3}{14.1}\selectfont 把中国产业与社区实践置于全球分工，兼顾欧美积累与全球南方潜力。}
\shade[left color=teal!10,right color=teal!2,rounded corners=2mm,draw=teal!17,line width=.45pt] \Pos{108}{246} rectangle \Pos{188}{273};
\path[fill=teal!9] \Pos{136.8}{273} .. controls \Pos{162.4}{272} and \Pos{168.8}{268} .. \Pos{187}{269} -- \Pos{187}{272} -- \Pos{136.8}{273} -- cycle;
\Box{112.00}{248.50}{72.00}{\sffamily\bfseries\fontsize{10.4}{15}\selectfont\color{teal}深度洞察 · AI 普惠}
\Box{112.00}{256.50}{72.00}{\fontsize{9.3}{14.1}\selectfont 从地域差异与时序演化识别战略机会，服务产业布局、地方决策与公共能力建设。}

\end{canvas}

% 06 — Four strategic insights
\begin{canvas}
\Header{执行摘要 · 关键洞察与主张}{teal}\Footer{06}
\Kicker{teal}{EXECUTIVE SUMMARY / INSIGHTS}\Title{执行摘要}
\Subhead{60}{blue}{关键洞察与主张}
\Box{22.00}{79.00}{14.00}{\digits\bfseries\fontsize{20}{25}\selectfont\color{blue}01}
\Box{40.00}{79.00}{148.00}{\sffamily\bfseries\fontsize{11.2}{16}\selectfont 全球协作的深度，塑造开放生态的体系能力}
\Box{40.00}{89.00}{148.00}{\fontsize{10.4}{16.4}\selectfont 欧洲的人才积累、中国的模型供给、印度与巴西的参与基础，体现全球开源的地域差异；增长速度与既有规模也不总是同向。将人才、组织、项目、代码、模型与数据统一观测，才能解释不同地区在创造、复用与维护中的分工及其时序变化。本报告主张把跨境接续、跨层支撑与长期维护纳入发展评价，以协作关系判断体系能力，为国际合作和公共投入提供依据。\textsuperscript{\color{teal}[1]}}
\Box{22.00}{125.00}{14.00}{\digits\bfseries\fontsize{20}{25}\selectfont\color{blue}02}
\Box{40.00}{125.00}{148.00}{\sffamily\bfseries\fontsize{11.2}{16}\selectfont 中国模型开源的领先，需要拓展为体系与协作优势}
\Box{40.00}{135.00}{148.00}{\fontsize{10.4}{16.4}\selectfont HF 期后观察中，Qwen 的衍生仓库达到 15.1 万个，显示以其为代表的中国模型处于全球开放复用生态的领先行列。围绕模型，训练推理、AI 工具链与异构算力适配正在分层展开；国家政策纳入开源贡献激励，xLLM 等项目进入基金会治理。下一步应以开放接口、共同评测与跨境维护连接国内产业和全球社区，使模型优势沉淀为长期体系能力。\textsuperscript{\color{teal}[2]}}
\Box{22.00}{171.00}{14.00}{\digits\bfseries\fontsize{20}{25}\selectfont\color{blue}03}
\Box{40.00}{171.00}{148.00}{\sffamily\bfseries\fontsize{11.2}{16}\selectfont 全球南方的潜力，在于成为共创者与治理参与者}
\Box{40.00}{181.00}{148.00}{\fontsize{10.4}{16.4}\selectfont 本地语言、文化知识与真实场景，是全球南方参与 AI 开源的独特资源。Masakhane 的非洲语言协作展示了以数据共建、模型适配和社区研究进入全球创新网络的路径。基金会应发挥连接长期资金、中立治理、人才培养与跨境合作的关键作用。AI 普惠应支持当地伙伴共同定义问题、掌握维护能力并参与规则制定，把资源可及推进为可持续公共能力。\textsuperscript{\color{teal}[3]}}
\Box{22.00}{217.00}{14.00}{\digits\bfseries\fontsize{20}{25}\selectfont\color{blue}04}
\Box{40.00}{217.00}{148.00}{\sffamily\bfseries\fontsize{11.2}{16}\selectfont 杭州等成长节点提示：中国 AI 开源需要因地施策}
\Box{40.00}{227.00}{148.00}{\fontsize{10.4}{16.4}\selectfont 中国 AI 开源呈现企业牵引、科研供给与社区接续并存的结构，模型、工具链与基础软件承担不同产业角色。人才规模、增长速度与模型供给指向不同的优势节点；杭州的开源人才增长与浙江的模型供给值得联读。地方决策应区分研究、工程适配、长期维护与应用人才需求，围绕自身产业补齐协作环节，通过跨城、跨境项目合作形成特色。\textsuperscript{\color{teal}[4]}}

\Box{22}{265}{166}{\sffamily\fontsize{7.6}{10.8}\selectfont\color{muted}证据入口：[1] 第 2—5 章的地图、榜单、技术全景与模型网络；[2] \href{\HF}{HF 夏季观察}（期后）、\href{https://www.cac.gov.cn/2025-08/27/c_1758018277755538.htm}{“人工智能+”意见}（2025-08）、\href{https://openatom.cn/journalism/detail/0uZHDVIiwISV}{开放原子项目捐赠}（2026-06）；[3] \href{https://www.masakhane.io/}{Masakhane} 与 \href{https://idrc-crdi.ca/en/what-we-do/projects-we-support/project/masakhane-ai-languages-hub}{IDRC 项目资料}（2026-11 查阅）；[4] 2.4 固定十城与 5.1 省级样本，各保留自身窗口，人才筛选细则待核，不推算全市 AI 人才总量。}
\end{canvas}

% 07 — 2026 declaration in the 2025 preface layout
\begin{canvas}
\hypertarget{intro}{}\pdfbookmark[0]{引言：全球开源与人工智能国际合作宣言}{intro-bm}
\Header{引言}{teal}\Footer{07}
\Kicker{teal}{INTRODUCTION / DECLARATION}\Title{引言}
\Box{22}{63}{166}{\sffamily\fontsize{9.3}{15.2}\selectfont\color{muted}本报告以《全球开源与人工智能国际合作宣言》为思想起点，以“开源共创智能未来”为共同愿景。通过跨平台数据、协作关系与实施案例，报告观察人才、组织、项目、代码、模型与数据如何共同演进，追问开放资源怎样转化为可持续公共能力。以下保留 2026 版宣言全文，并以结构图呈现其愿景、原则与行动共识。}
\Box{22}{100}{166}{\centering\sffamily\bfseries\fontsize{16}{23}\selectfont\color{teal}《全球开源与人工智能国际合作宣言》}
\Box{22}{112}{166}{\centering\sffamily\fontsize{10.5}{16}\selectfont\color{muted}—— 开源共创智能未来：迈向全球 AI 治理新纪元}
\Box{22}{128}{166}{\DeclBody\DHead{I. 时代呼唤：迈向智能协作新纪元}
我们正身处人类文明从“计算互联”向“智能共生”跃迁的关键节点。人工智能技术的爆发式增长，特别是大模型与 AI 智能体的崛起，以前所未有的深度重塑着全球经济与社会结构。然而，全球智能鸿沟持续扩大，技术主权博弈加剧，数据孤岛与算法黑盒成为悬在全人类头顶的达摩克利斯之剑。面对气候变化、公共卫生挑战以及复杂的国际关系，唯有坚持开源共享、协作共治，才能确保 AI 向善，避免技术垄断带来的不平等。

在此，我们 —— 来自中国政府、世界顶尖科技企业、国际学术机构、开源基金会及全球开发者社区的代表——郑重发布本宣言。我们呼吁以开放精神为指引，依托世界人工智能合作组织（WAICO）等多边平台，共建共享一个更加公平、透明、可持续的智能未来。

本宣言提出全球 AI 开源行动框架：以数字基础设施为基石，以实质贡献与伦理治理为支柱，以社会价值实现为使命，推动开源软件成为驱动智能时代可持续发展的核心引擎。

\DHead{II. 共同愿景：构建智能共治的数字文明}
我们致力于构建一个开放、包容、安全、可持续的全球 AI 开源生态，使强大的智能技术真正服务于全人类福祉。

我们坚信：开源不仅是代码的共享，更是信任的建立。在生成式 AI 与大模型时代，开源构成了全球数字公共品的底层架构。唯有坚持多元参与、透明可解释、贡献导向的协作原则，才能让技术创新成为推动和平、繁荣与可持续发展的共同力量。

我们共同倡导：将开源创新纳入全球数字公共品战略，作为实现《2030 年可持续发展议程》的关键路径，助力构建共治、共享、共赢的数字文明。
}
\end{canvas}

% 08 — Original declaration, second part
\begin{canvas}
\Header{引言 · 宣言全文（续）}{teal}\Footer{08}
\Kicker{teal}{INTRODUCTION / DECLARATION}\Title{引言\hspace{3mm}{\fontsize{13}{18}\selectfont（续）}}
\Box{22}{63}{166}{\DeclBody\DHead{III. 核心原则：行动框架}
\DLayer{第一层：筑牢智能时代基石}
\DPrinciple{原则一：坚持开放互联}{“万物智联，普惠共生”}
打破算法壁垒与语言障碍，推动多模态模型的跨文化交互。鼓励支持多语言的本地化大模型开发，坚持技术发展的非歧视性。让发展中国家能够平等接入全球 AI 资源，使技术连接成为缩小数字鸿沟的桥梁，而非加剧不平等的工具。

\DPrinciple{原则二：坚持安全可信}{“透明算法，信任基石”}
反对互联网碎片化与技术封闭。支持开源作为验证算法安全、检测数据偏见的首选试验田。建立可信、弹性、可持续的 AI 供应链安全机制，确保关键基础设施的中立性与稳定性，为全球智能社会的运行筑牢底线。

\DLayer{第二层：夯实伦理治理支柱}
\DPrinciple{原则三：坚持数据伦理}{“算法向善，责任共担”}
坚持数据隐私保护与算法公平性，防范技术滥用。推动 AI 治理遵循包容、敏捷、负责任的原则。鼓励开源社区探索可解释性 AI（XAI）的最佳实践，确保每一项 AI 决策都可追溯、可审计，始终服务于人类尊严与权利。

\DPrinciple{原则四：坚持多元包容}{“众智共创，价值共享”}
尊重文化多样性与认知能力的差异，反对任何形式的“算法歧视”。特别关注弱势群体在 AI 时代的权益，通过开源项目赋能边缘群体。支持多元化的开发者社群，让每一个声音都能在智能技术的演进中被听见、被尊重。

\DPrinciple{原则五：坚持全球协作}{“跨界协同，智能互联”}
视开源为连接全球智慧的网络。支持跨国界、跨机构的联合研发，特别是在应对气候变化、生物医药等全球性挑战领域。促进知识、算力与高质量数据集的自由流动，强化南南合作与多边主义，构建去中心化但互联互通的创新生态。

\DLayer{第三层：共创社会价值未来}
\DPrinciple{原则六：坚持绿色智能}{“低碳驱动，创新致远”}
优先支持应对气候变化、环境保护、清洁能源等议题的开源 AI 项目。倡导高效的模型训练范式与算力调度优化，降低 AI 运行的碳足迹。推动开源技术作为数字公共品的核心实现路径，在能源管理、环境监测中的广泛应用，确保智能技术的发展不以牺牲生态环境为代价。
}
\end{canvas}
\begin{canvas}
\Header{引言 · 宣言全文（续）}{teal}\Footer{09}
\Kicker{teal}{INTRODUCTION / DECLARATION}\Title{引言\hspace{3mm}{\fontsize{13}{18}\selectfont（续）}}
\Box{22}{63}{166}{\DeclBody\DPrinciple{原则七：坚持开源普惠}{“智能为公，成果共享”}
推动开源技术广泛应用于政府治理、公共卫生、教育公平等领域，服务社会治理现代化与民生改善。鼓励开源成为数字公共品的载体，建立以贡献为导向的激励机制，让创新成果惠及所有人，尤其是边缘群体。

\DHead{IV. 行动共识：迈向全球 AI 治理新格局}
我们呼吁：各国政府、科技企业、学界与公民社会携手同行，以本宣言为基础，深化在 AI 治理、算力共享、人才培育等领域的务实合作，争取早期收获。

我们支持依托联合国及相关国际机制，围绕人工智能伦理准则、跨境数据流动规则、AI 版权保护等紧迫议题展开对话，凝聚共识，推动开源成为构建公正合理全球人工智能治理体系的生动实践。

\DHead{结语}
开源，是我们对开放精神的坚守，是对人类团结的信念，是对更美好世界的共同追求。让我们以协作代替割裂，以透明回应质疑，以共享消弭垄断。共同构建一个更加开放、公平、可持续的智能文明。

开源创新，属于全人类；智能未来，由我们共同创造。
}
\Box{22}{181}{166}{\sffamily\bfseries\fontsize{10.3}{15}\selectfont\color{navy}宣言总览｜以共同愿景贯通三层原则}
\draw[blue!10,line width=4pt] \Pos{64}{224} .. controls \Pos{72}{224} and \Pos{72}{199.5} .. \Pos{79}{199.5};
\draw[blue!55,line width=0.65pt] \Pos{64}{224} .. controls \Pos{72}{224} and \Pos{72}{199.5} .. \Pos{79}{199.5};
\draw[teal!10,line width=4pt] \Pos{64}{224} .. controls \Pos{72}{224} and \Pos{72}{223.5} .. \Pos{79}{223.5};
\draw[teal!55,line width=0.65pt] \Pos{64}{224} .. controls \Pos{72}{224} and \Pos{72}{223.5} .. \Pos{79}{223.5};
\draw[violet!10,line width=4pt] \Pos{64}{224} .. controls \Pos{72}{224} and \Pos{72}{247.5} .. \Pos{79}{247.5};
\draw[violet!55,line width=0.65pt] \Pos{64}{224} .. controls \Pos{72}{224} and \Pos{72}{247.5} .. \Pos{79}{247.5};
\fill[blue!5] \Pos{46}{224} circle[radius=18.2mm];
\shade[top color=white,bottom color=blue!28,shading angle=30,draw=blue!35,line width=.45pt] \Pos{46}{224} circle[radius=17mm];
\draw[white,line width=.7pt] \Pos{36.14}{214.14} arc[start angle=135,end angle=65,radius=13.94mm];
\Box{27}{199}{38}{\sffamily\bfseries\fontsize{8.5}{12}\selectfont\color{teal}\centering 共同愿景}
\Box{28}{216}{36}{\sffamily\bfseries\fontsize{12.2}{17}\selectfont\color{navy}\centering 开源共创}
\Box{28}{227}{36}{\sffamily\bfseries\fontsize{12.2}{17}\selectfont\color{navy}\centering 智能未来}
\shade[left color=blue!10,right color=blue!2,rounded corners=2mm,draw=blue!17,line width=.45pt] \Pos{80}{190} rectangle \Pos{188}{209};
\path[fill=blue!9] \Pos{118.88}{209} .. controls \Pos{153.44}{208} and \Pos{162.07999999999998}{204} .. \Pos{187}{205} -- \Pos{187}{208} -- \Pos{118.88}{209} -- cycle;
\Box{84}{192}{100}{\sffamily\bfseries\fontsize{8.7}{12}\selectfont\color{blue}智能时代基石}
\Box{84.0}{200}{4.1}{\sffamily\fontsize{9.5}{14}\selectfont\color{ink}\includegraphics[width=4.1mm]{.build/assets/principle-1.pdf}}
\Box{90.0}{200}{43.0}{\sffamily\bfseries\fontsize{9.2}{12}\selectfont\color{navy}开放互联}
\Box{134.0}{200}{4.1}{\sffamily\fontsize{9.5}{14}\selectfont\color{ink}\includegraphics[width=4.1mm]{.build/assets/principle-2.pdf}}
\Box{140.0}{200}{43.0}{\sffamily\bfseries\fontsize{9.2}{12}\selectfont\color{navy}安全可信}
\shade[left color=teal!10,right color=teal!2,rounded corners=2mm,draw=teal!17,line width=.45pt] \Pos{80}{214} rectangle \Pos{188}{233};
\path[fill=teal!9] \Pos{118.88}{233} .. controls \Pos{153.44}{232} and \Pos{162.07999999999998}{228} .. \Pos{187}{229} -- \Pos{187}{232} -- \Pos{118.88}{233} -- cycle;
\Box{84}{216}{100}{\sffamily\bfseries\fontsize{8.7}{12}\selectfont\color{teal}伦理治理支柱}
\Box{84.0}{224}{4.1}{\sffamily\fontsize{9.5}{14}\selectfont\color{ink}\includegraphics[width=4.1mm]{.build/assets/principle-3.pdf}}
\Box{90.0}{224}{26.333333333333336}{\sffamily\bfseries\fontsize{9.2}{12}\selectfont\color{navy}数据伦理}
\Box{117.33333333333334}{224}{4.1}{\sffamily\fontsize{9.5}{14}\selectfont\color{ink}\includegraphics[width=4.1mm]{.build/assets/principle-4.pdf}}
\Box{123.33333333333334}{224}{26.333333333333336}{\sffamily\bfseries\fontsize{9.2}{12}\selectfont\color{navy}多元包容}
\Box{150.66666666666669}{224}{4.1}{\sffamily\fontsize{9.5}{14}\selectfont\color{ink}\includegraphics[width=4.1mm]{.build/assets/principle-5.pdf}}
\Box{156.66666666666669}{224}{26.333333333333336}{\sffamily\bfseries\fontsize{9.2}{12}\selectfont\color{navy}全球协作}
\shade[left color=violet!10,right color=violet!2,rounded corners=2mm,draw=violet!17,line width=.45pt] \Pos{80}{238} rectangle \Pos{188}{257};
\path[fill=violet!9] \Pos{118.88}{257} .. controls \Pos{153.44}{256} and \Pos{162.07999999999998}{252} .. \Pos{187}{253} -- \Pos{187}{256} -- \Pos{118.88}{257} -- cycle;
\Box{84}{240}{100}{\sffamily\bfseries\fontsize{8.7}{12}\selectfont\color{violet}社会价值未来}
\Box{84.0}{248}{4.1}{\sffamily\fontsize{9.5}{14}\selectfont\color{ink}\includegraphics[width=4.1mm]{.build/assets/principle-6.pdf}}
\Box{90.0}{248}{43.0}{\sffamily\bfseries\fontsize{9.2}{12}\selectfont\color{navy}绿色智能}
\Box{134.0}{248}{4.1}{\sffamily\fontsize{9.5}{14}\selectfont\color{ink}\includegraphics[width=4.1mm]{.build/assets/principle-7.pdf}}
\Box{140.0}{248}{43.0}{\sffamily\bfseries\fontsize{9.2}{12}\selectfont\color{navy}开源普惠}
\shade[left color=teal!10,right color=teal!2,rounded corners=2mm,draw=teal!17,line width=.45pt] \Pos{22}{261} rectangle \Pos{188}{270};
\path[fill=teal!9] \Pos{81.75999999999999}{270} .. controls \Pos{134.88}{269} and \Pos{148.16}{265} .. \Pos{187}{266} -- \Pos{187}{269} -- \Pos{81.75999999999999}{270} -- cycle;
\Box{27}{262.5}{28}{\sffamily\bfseries\fontsize{9.5}{13}\selectfont\color{teal}行动共识}
\Box{60}{262.5}{123}{\sffamily\bfseries\fontsize{9.4}{13}\selectfont\color{navy}深化 AI 治理、算力共享与人才培育合作}
\Box{22}{273}{166}{\sffamily\fontsize{7.8}{11}\selectfont\color{muted}图引 1　2026 宣言的共同愿景、三层七项原则与行动共识（依据前述全文绘制）。}

\end{canvas}
\end{document}

